\documentclass[aspectratio=169,11pt]{beamer}
\usepackage[T1]{fontenc}
\usepackage[utf8]{inputenc}
\usepackage[spanish]{babel}
\uselanguage{Spanish}
\languagepath{Spanish}
\usepackage{lmodern}
\usepackage{amsmath,amssymb,mathtools}
\usepackage{xcolor}
\usepackage{tikz}
\usetikzlibrary{calc,positioning}

\definecolor{UdeAGreen}{HTML}{0D6939}
\definecolor{DeepGreen}{HTML}{064629}
\definecolor{LeafGreen}{HTML}{83B840}
\definecolor{PaleGreen}{HTML}{E7F0E7}
\definecolor{CardGreen}{HTML}{F5F8F3}
\definecolor{Ink}{HTML}{173629}
\definecolor{WarmGold}{HTML}{E5C45A}

\usetheme{Boadilla}
\setbeamertemplate{navigation symbols}{}
\setbeamertemplate{blocks}[rounded][shadow=false]
\setbeamercolor{normal text}{fg=Ink,bg=PaleGreen}
\setbeamercolor{background canvas}{bg=PaleGreen}
\setbeamercolor{structure}{fg=UdeAGreen}
\setbeamercolor{frametitle}{fg=white,bg=DeepGreen}
\setbeamerfont{frametitle}{series=\bfseries,size=\large}
\setbeamercolor{title}{fg=white}
\setbeamercolor{subtitle}{fg=white!85}
\setbeamercolor{block title}{fg=white,bg=UdeAGreen}
\setbeamercolor{block body}{fg=Ink,bg=CardGreen}
\setbeamercolor{block title example}{fg=DeepGreen,bg=LeafGreen!55}
\setbeamercolor{block body example}{fg=Ink,bg=CardGreen}
\setbeamercolor{block title alerted}{fg=white,bg=DeepGreen}
\setbeamercolor{block body alerted}{fg=Ink,bg=CardGreen}
\setbeamercolor{author in head/foot}{fg=white,bg=DeepGreen}
\setbeamertemplate{footline}{%
  \leavevmode%
  \hbox{\begin{beamercolorbox}[wd=\paperwidth,ht=2.6ex,dp=1.1ex,leftskip=1em,rightskip=1em]{author in head/foot}%
  \insertshorttitle\hfill Grupo de Estudio · UdeA\hfill\insertframenumber/\inserttotalframenumber
  \end{beamercolorbox}}%
}
\setbeamertemplate{itemize item}{\color{UdeAGreen}\small$\blacktriangleright$}
\setbeamertemplate{itemize subitem}{\color{LeafGreen}\scriptsize$\blacktriangleright$}
\setlength{\leftmargini}{1.5em}

\newcommand{\NN}{\mathbb{N}}
\newcommand{\ZZ}{\mathbb{Z}}
\newcommand{\RR}{\mathbb{R}}
\newcommand{\PP}{\mathcal{P}}
\newcommand{\card}[1]{\lvert #1\rvert}
\newcommand{\Set}[1]{\left\{#1\right\}}

\title[Teoría de conjuntos]{Teoría de conjuntos}
\subtitle{Clase 1 · Lenguaje, operaciones y demostraciones}
\author{Daniel Soto Villada}
\institute{Grupo de Estudio de Teoría Analítica de Números\\Facultad de Ciencias Exactas y Naturales · Universidad de Antioquia}
\date{12 de noviembre de 2026}

\begin{document}

{\setbeamercolor{background canvas}{bg=DeepGreen}
\begin{frame}[plain]
    \begin{tikzpicture}[remember picture,overlay]
        \node[anchor=center] at ([xshift=-2cm]current page.east)
            {\includegraphics[height=6.1cm,keepaspectratio]{LogoUdeA_Blanco.png}};
    \end{tikzpicture}
    \vspace{0.65cm}
    {\color{WarmGold}\bfseries\large GRUPO DE ESTUDIO}\par
    \vspace{0.35cm}
    {\color{white}\bfseries\fontsize{29}{34}\selectfont TEORÍA ANALÍTICA\\DE NÚMEROS}\par
    \vspace{0.55cm}
    {\color{LeafGreen}\bfseries\large INICIO · PRIMERA UNIDAD}\par
    \vspace{0.12cm}
    {\color{white}\Large Teoría de conjuntos}\par
    {\color{white!80}\normalsize Lenguaje, operaciones y demostraciones}\par
    \vspace{0.75cm}
    {\color{white}\large Daniel Soto Villada}\par
    \vspace{0.15cm}
    {\color{white!80}\small Facultad de Ciencias Exactas y Naturales · Universidad de Antioquia}\par
    \vfill
\end{frame}}

\section{Conjuntos: notación y propiedades}
\begin{frame}{Conjuntos}
    \begin{definition}
        Un conjunto es una colección bien definida de objetos distintos, llamados elementos. Si un objeto $x$ pertenece al conjunto $A$, escribimos $x\in A$; si no pertenece, escribimos $x\notin A$.
    \end{definition}
    \begin{columns}[T,onlytextwidth]
        \column{.47\textwidth}
        \begin{exampleblock}{Ejemplo}
            Para $A=\Set{a,b,c}$, se tiene $a\in A$ y $d\notin A$.
        \end{exampleblock}
        \begin{block}{Criterio de igualdad}
            Dos conjuntos son iguales cuando tienen los mismos elementos; el orden y las repeticiones no los alteran.
        \end{block}
        \column{.49\textwidth}
        \centering
    \begin{tikzpicture}[scale=.9]
    \draw[DeepGreen,thick,rounded corners] (0,0) rectangle (4.5,2.5);
    \node[anchor=north west,text=DeepGreen] at (.2,2.35) {$A$};
    \fill[UdeAGreen] (1,1.5) circle (2pt);
    \node[above] at (1,1.5) {$a$};
    \fill[UdeAGreen] (2.25,.8) circle (2pt);
    \node[below] at (2.25,.8) {$b$};
    \fill[UdeAGreen] (3.5,1.5) circle (2pt);
    \node[above] at (3.5,1.5) {$c$};
    \fill[LeafGreen] (5.1,1.25) circle (2pt);
    \node[right=2pt] at (5.1,1.25) {$d$};
\end{tikzpicture}
        \small Los puntos interiores representan los elementos de $A$.
    \end{columns}
\end{frame}

\begin{frame}{Pertenencia}
    \begin{block}{Definición}
        Si $A$ es un conjunto y $x$ es uno de sus elementos, escribimos $x\in A$. Si no es elemento, escribimos $x\notin A$.
    \end{block}
    \begin{exampleblock}{Ejemplo}
        Para $A=\Set{\text{rojo},\text{verde},\text{azul}}$,
        \[
          \text{verde}\in A,\qquad \text{amarillo}\notin A.
        \]
    \end{exampleblock}
    \begin{alertblock}{Distinción de notación}
        $\in$ relaciona un elemento con un conjunto. No significa lo mismo que $\subseteq$, que compara dos conjuntos.
    \end{alertblock}
\end{frame}

\begin{frame}{Igualdad de conjuntos}
    \begin{block}{Enunciado}
        Sean
        \[
          A=\Set{x\in\ZZ:0\leq x\leq4\text{ y }x\text{ es par}},
          \qquad B=\Set{0,2,4}.
        \]
        Entonces $A=B$.
    \end{block}
    \begin{block}{Demostración}
        Si $x\in A$, entonces $x$ es un entero par entre $0$ y $4$; por tanto, $x\in\Set{0,2,4}=B$, y $A\subseteq B$.

        Recíprocamente, cada elemento de $B$ es un entero par entre $0$ y $4$, así que pertenece a $A$; luego $B\subseteq A$. Como ambos conjuntos tienen los mismos elementos, $A=B$.
    \end{block}
\end{frame}


\begin{frame}{Extensionalidad e identidad de los elementos}
    \begin{columns}[T,onlytextwidth]
        \column{.52\textwidth}
        \begin{exampleblock}{Mismos elementos}
            \[
            \Set{1,2,2,3}=\Set{3,1,2}=\Set{1,2,3}.
            \]
            El orden no importa y repetir un elemento no crea uno nuevo.
        \end{exampleblock}
        \column{.44\textwidth}
        \begin{alertblock}{Pero el tipo de objeto sí importa}
            \[
            1\ne\Set{1},\qquad \Set{1}\ne\Set{\Set{1}}.
            \]
            El elemento $1$ no es el conjunto cuyo único elemento es $1$.
        \end{alertblock}
    \end{columns}
    \begin{exampleblock}{Tarea de programación}
        \href{https://codeforces.com/problemset/problem/443/A}{\textit{Anton and Letters} (Codeforces)}.
    \end{exampleblock}
\end{frame}

\begin{frame}{Representación de conjuntos}
    \begin{columns}[T,onlytextwidth]
        \column{.47\textwidth}
        \begin{block}{Por extensión}
            Se enumeran sus elementos:
            \[
              A=\Set{2,4,6,8}.
            \]
        \end{block}
        \column{.49\textwidth}
        \begin{block}{Por comprensión}
            Se indica una propiedad que cumplen:
            \[
              A=\Set{x\in\NN: x\text{ es par y }0<x<10}.
            \]
        \end{block}
    \end{columns}
    \begin{exampleblock}{Ejemplo infinito}
        \[
          P=\Set{n\in\NN:n\text{ es primo}}=\Set{2,3,5,7,11,\ldots}.
        \]
    \end{exampleblock}
    \small En la notación $\Set{x\in X:P(x)}$, se consideran los elementos de $X$ que satisfacen $P$.
\end{frame}

\begin{frame}{Conjunto universo y comprensión}
    \begin{block}{Conjunto definido dentro de un universo}
        Para un conjunto $X$ y una propiedad $P$, escribimos
        \[
          A=\Set{x\in X:P(x)}.
        \]
        El universo $X$ especifica de dónde pueden venir los elementos.
    \end{block}
    \begin{exampleblock}{La misma propiedad, distintos universos}
        \[
        \Set{x\in\ZZ:x^2=1}=\Set{-1,1},\qquad
        \Set{x\in\NN:x^2=1}=\Set{1},
        \]
        usando aquí $\NN=\Set{0,1,2,\ldots}$.
    \end{exampleblock}
    \begin{exampleblock}{Tarea de programación}
        \href{https://codeforces.com/problemset/problem/520/A}{\textit{Pangram} (Codeforces)}.
    \end{exampleblock}
\end{frame}

\begin{frame}{Inclusión de conjuntos}
    \begin{block}{Definición}
        $A$ es subconjunto de $B$, y escribimos $A\subseteq B$, cuando
        \[
          \forall x\;(x\in A\Longrightarrow x\in B).
        \]
    \end{block}
    \begin{exampleblock}{Ejemplo numérico}
        Si $A=\Set{2,4}$ y $B=\Set{1,2,3,4,5}$, entonces $A\subseteq B$.
    \end{exampleblock}
    \begin{alertblock}{No confundir}
        $2\in B$ es una afirmación de pertenencia; $\Set{2}\subseteq B$ es una afirmación de inclusión.
    \end{alertblock}
\end{frame}

\begin{frame}{Diagrama de inclusión}
    \centering
    \begin{tikzpicture}[scale=.95]
    \draw[DeepGreen,thick,rounded corners] (0,0) rectangle (6.5,3.5);
    \node[anchor=north east,text=DeepGreen] at (6.3,3.3) {$B$};
    \draw[UdeAGreen,thick] (2.3,1.65) ellipse (1.7 and 1.15);
    \node[above left,text=UdeAGreen] at (1,2.25) {$A$};
    \fill[UdeAGreen] (1.8,1.55) circle (2pt);
    \node[below left] at (1.8,1.55) {$a$};
    \fill[UdeAGreen] (2.9,1.65) circle (2pt);
    \node[below right] at (2.9,1.65) {$b$};
    \fill[LeafGreen] (5.3,1.2) circle (2pt);
    \node[below] at (5.3,1.2) {$c$};
\end{tikzpicture}
    \begin{exampleblock}{Lectura del diagrama}
        Si $A=\Set{a,b}$ y $B=\Set{a,b,c}$, entonces $A\subseteq B$.
    \end{exampleblock}
\end{frame}

\begin{frame}{Demostración de una inclusión}
    \begin{block}{Enunciado}
        Sean $A=\Set{4k:k\in\ZZ}$, el conjunto de múltiplos de $4$, y $B=\Set{2m:m\in\ZZ}$, el conjunto de enteros pares. Probaremos que $A\subseteq B$.
    \end{block}
    \begin{block}{Demostración}
        Tomemos un elemento arbitrario $x\in A$. Por definición de $A$, existe $k\in\ZZ$ tal que $x=4k$. Definimos $m=2k$; como $k$ es entero, $m\in\ZZ$. Entonces
        \[
          x=4k=2(2k)=2m.
        \]
        Por definición de $B$, esto implica que $x\in B$. Como $x$ fue arbitrario, concluimos que $A\subseteq B$.
    \end{block}
    \begin{exampleblock}{Caso concreto}
        $12\in A$ porque $12=4\cdot3$; tomando $m=2\cdot3=6$, también $12=2\cdot6$, así que $12\in B$.
    \end{exampleblock}
\end{frame}

\begin{frame}{Criterio de igualdad de conjuntos}
    \begin{theorem}[Criterio de doble inclusión]
        Para cualesquiera conjuntos $A$ y $B$,
        \[
          A=B\quad\Longleftrightarrow\quad A\subseteq B\text{ y }B\subseteq A.
        \]
    \end{theorem}
    \begin{block}{Demostración}
        Si $A=B$, tienen los mismos elementos, así que ambas inclusiones se cumplen.

        Recíprocamente, supongamos $A\subseteq B$ y $B\subseteq A$. Para todo $x$, la primera inclusión da $x\in A\Rightarrow x\in B$, y la segunda da $x\in B\Rightarrow x\in A$. Así, $x\in A$ si y solo si $x\in B$; por extensionalidad, $A=B$.
    \end{block}
\end{frame}

\begin{frame}{Ejemplo de doble inclusión}
    Sean
    \[
      A=\Set{x\in\ZZ:x\text{ es par y }0\leq x\leq 8},\qquad
      B=\Set{0,2,4,6,8}.
    \]
    \begin{enumerate}
        \item Si $x\in A$, es un entero par entre $0$ y $8$; por tanto $x$ es uno de $0,2,4,6,8$. Luego $A\subseteq B$.
        \item Si $x\in B$, es uno de $0,2,4,6,8$; cada uno es par y está entre $0$ y $8$. Luego $B\subseteq A$.
    \end{enumerate}
    \begin{exampleblock}{Conclusión}
        Por el criterio de doble inclusión, $A=B$.
    \end{exampleblock}

\end{frame}

\begin{frame}{Conjunto vacío y conjuntos unitarios}
    \begin{block}{Definiciones}
        El conjunto sin elementos se denota $\varnothing$. Un conjunto con exactamente un elemento, como $\Set{a}$, se llama unitario.
    \end{block}
    \begin{theorem}
        Para todo conjunto $A$, se cumple $\varnothing\subseteq A$.
    \end{theorem}
    \begin{block}{Demostración}
        Si $\varnothing\nsubseteq A$, existiría un elemento $x\in\varnothing$ con $x\notin A$. Pero el conjunto vacío no tiene elementos. La suposición es imposible; por tanto $\varnothing\subseteq A$.
    \end{block}
    \begin{exampleblock}{Pertenencia e inclusión}
        $\varnothing\subseteq\Set{a}$ siempre; en cambio, $\varnothing\in\Set{a}$ solo si $a=\varnothing$.
    \end{exampleblock}
\end{frame}

\begin{frame}{Conjunto potencia}
    \begin{definition}
        El conjunto potencia de $A$, denotado $\PP(A)$, es el conjunto de todos los subconjuntos de $A$; también se llama conjunto de partes de $A$:
        \[
          \PP(A)=\Set{S:S\subseteq A}.
        \]
    \end{definition}
    \begin{exampleblock}{Ejemplo}
        Si $A=\Set{a,b}$, entonces
        \[
          \PP(A)=\Set{\varnothing,\Set{a},\Set{b},\Set{a,b}}.
        \]
    \end{exampleblock}
    \begin{alertblock}{Pertenencia e inclusión}
        $a\in A$, pero $\Set{a}\subseteq A$ y $\Set{a}\in\PP(A)$. Son afirmaciones distintas.
    \end{alertblock}
\end{frame}

\begin{frame}{Ejemplo: conjunto potencia de tres elementos}
    \centering
    \begin{tikzpicture}[x=1cm,y=.76cm,>=stealth,setnode/.style={draw=DeepGreen,thick,rounded corners,fill=CardGreen,minimum width=1.65cm,minimum height=.56cm,inner sep=3pt}]
    \node[setnode] (empty) at (0,-2.5) {$\varnothing$};
    \node[setnode] (a) at (2.5,0) {$\Set{a}$};
    \node[setnode] (ab) at (0,2.5) {$\Set{a,b}$};
    \node[setnode] (b) at (-2.5,0) {$\Set{b}$};
    \node[setnode] (c) at (0,.2) {$\Set{c}$};
    \node[setnode] (ac) at (1.9,1.9) {$\Set{a,c}$};
    \node[setnode] (abc) at (0,3.6) {$\Set{a,b,c}$};
    \node[setnode] (bc) at (-1.9,1.9) {$\Set{b,c}$};

    \draw[->,thick,UdeAGreen] (empty) -- (a);
    \draw[->,thick,UdeAGreen] (empty) -- (b);
    \draw[->,thick,UdeAGreen] (a) -- (ab);
    \draw[->,thick,UdeAGreen] (b) -- (ab);
    \draw[->,thick,UdeAGreen] (c) -- (ac);
    \draw[->,thick,UdeAGreen] (c) -- (bc);
    \draw[->,thick,UdeAGreen] (ac) -- (abc);
    \draw[->,thick,UdeAGreen] (bc) -- (abc);
    \draw[->,thick,UdeAGreen] (empty) -- (c);
    \draw[->,thick,UdeAGreen] (a) -- (ac);
    \draw[->,thick,UdeAGreen] (b) -- (bc);
    \draw[->,thick,UdeAGreen] (ab) -- (abc);
\end{tikzpicture}
    \begin{block}{Lectura del diagrama}
        Para $A=\Set{a,b,c}$, el diagrama muestra sus ocho subconjuntos: $\card{\PP(A)}=2^3=8$. Cada flecha indica una inclusión inmediata.
    \end{block}
\end{frame}

\section{Operaciones y propiedades}
\begin{frame}{Unión e intersección}
    \begin{columns}[T,onlytextwidth]
        \column{.49\textwidth}
        \begin{block}{Unión}
            \[
              A\cup B=\Set{x:x\in A\text{ o }x\in B}.
            \]
            El «o» es inclusivo: puede pertenecer a ambos.
        \end{block}
        \begin{exampleblock}{Ejemplo}
            $\Set{1,2,3}\cup\Set{3,4}=\Set{1,2,3,4}$.
        \end{exampleblock}
        \column{.49\textwidth}
        \begin{block}{Intersección}
            \[
              A\cap B=\Set{x:x\in A\text{ y }x\in B}.
            \]
            Se conservan los elementos comunes.
        \end{block}
        \begin{exampleblock}{Ejemplo}
            $\Set{1,2,3}\cap\Set{3,4}=\Set{3}$.
        \end{exampleblock}
    \end{columns}
\end{frame}

\begin{frame}{Representación de unión e intersección}
    \centering
    \begin{tikzpicture}[scale=1.3]
    \begin{scope}[shift={(-3.1,0)}]
        \fill[CardGreen] (-1.65,-1.35) rectangle (1.65,1.35);
        \fill[LeafGreen!45] (-.55,0) circle (1.0);
        \fill[LeafGreen!45] (.55,0) circle (1.0);
        \draw[DeepGreen,thick,rounded corners] (-1.65,-1.35) rectangle (1.65,1.35);
        \draw[UdeAGreen,thick] (-.55,0) circle (1.0);
        \draw[UdeAGreen,thick] (.55,0) circle (1.0);
        \node at (-.9,.8) {$A$};\node at (.9,.8) {$B$};
        \node[anchor=north west,font=\small] at (-1.5,1.2) {$X$};
        \node[below=1.85cm,font=\small] at (0,0) {$A\cup B$};
    \end{scope}
    \begin{scope}[shift={(3.1,0)}]
        \fill[CardGreen] (-1.65,-1.35) rectangle (1.65,1.35);
        \begin{scope}
            \clip (-.55,0) circle (1.0);
            \fill[LeafGreen!65] (.55,0) circle (1.0);
        \end{scope}
        \draw[DeepGreen,thick,rounded corners] (-1.65,-1.35) rectangle (1.65,1.35);
        \draw[UdeAGreen,thick] (-.55,0) circle (1.0);
        \draw[UdeAGreen,thick] (.55,0) circle (1.0);
        \node at (-.9,.8) {$A$};\node at (.9,.8) {$B$};
        \node[anchor=north west,font=\small] at (-1.5,1.2) {$X$};
        \node[below=1.85cm,font=\small] at (0,0) {$A\cap B$};
    \end{scope}
\end{tikzpicture}\par
    \vspace{.1cm}
    \small El sombreado muestra la operación dentro del universo $X$.
\end{frame}

\begin{frame}{Diferencia y complemento}
    \begin{columns}[T,onlytextwidth]
        \column{.49\textwidth}
        \begin{block}{Diferencia}
            \[
              A\setminus B=\Set{x\in A:x\notin B}.
            \]
            \begin{exampleblock}{Ejemplo}
                $\Set{1,2,3,4}\setminus\Set{3,4,5}=\Set{1,2}$.
            \end{exampleblock}
        \end{block}
        \column{.49\textwidth}
        \begin{block}{Complemento}
            Fijado un universo $X$ y $A\subseteq X$,
            \[
              A^c=X\setminus A.
            \]
            \begin{exampleblock}{Ejemplo}
                Si $X=\Set{1,2,3,4,5}$ y $A=\Set{1,3,5}$, entonces $A^c=\Set{2,4}$.
            \end{exampleblock}
        \end{block}
    \end{columns}
\end{frame}

\begin{frame}{Ejemplo de cálculo de operaciones}
    Sean $A=\Set{1,2,3,4}$ y $B=\Set{3,4,5}$, con universo $X=\Set{1,2,3,4,5,6}$.
    \vspace{.2cm}
    \[
    \begin{array}{c|l}
        \text{Operación} & \text{Resultado}\\\hline
        A\cup B & \Set{1,2,3,4,5}\\
        A\cap B & \Set{3,4}\\
        A\setminus B & \Set{1,2}\\
        B\setminus A & \Set{5}\\
        A^c=X\setminus A & \Set{5,6}
    \end{array}
    \]
    \begin{alertblock}{Observación}
        $A\setminus B$ y $B\setminus A$ no tienen por qué ser iguales.
    \end{alertblock}
\end{frame}

\begin{frame}{Inclusiones básicas}
    \begin{theorem}
        Para cualesquiera $A$ y $B$,
        \[
          A\cap B\subseteq A\subseteq A\cup B.
        \]
    \end{theorem}
    \begin{block}{Demostración, elemento por elemento}
        Si $x\in A\cap B$, entonces $x\in A$ y $x\in B$; en particular, $x\in A$. Así $A\cap B\subseteq A$.

        Si $x\in A$, entonces cumple «$x\in A$ o $x\in B$»; por tanto $x\in A\cup B$. Así $A\subseteq A\cup B$.
    \end{block}
\end{frame}

\begin{frame}{Identidad distributiva}
    \begin{theorem}[Distributividad]
        \[
          A\cap(B\cup C)=(A\cap B)\cup(A\cap C).
        \]
    \end{theorem}
    \begin{block}{Demostración por pertenencia}
        Para cualquier $x$,
        \begin{align*}
          x\in A\cap(B\cup C)
          &\Longleftrightarrow x\in A\text{ y }(x\in B\text{ o }x\in C)\\
          &\Longleftrightarrow (x\in A\text{ y }x\in B)\text{ o }(x\in A\text{ y }x\in C)\\
          &\Longleftrightarrow x\in(A\cap B)\cup(A\cap C).
        \end{align*}
        Los dos conjuntos tienen exactamente los mismos elementos; son iguales.
    \end{block}
\end{frame}

\begin{frame}{Leyes de De Morgan}
    \begin{theorem}
        Para $A,B\subseteq X$,
        \[
          (A\cup B)^c=A^c\cap B^c,\qquad
          (A\cap B)^c=A^c\cup B^c.
        \]
    \end{theorem}
    \begin{block}{Demostración de la primera identidad}
        Para $x\in X$,
        \begin{align*}
          x\in(A\cup B)^c
          &\Longleftrightarrow x\notin A\cup B\\
          &\Longleftrightarrow (x\notin A)\text{ y }(x\notin B)\\
          &\Longleftrightarrow x\in A^c\cap B^c.
        \end{align*}
        La segunda se prueba de manera análoga, intercambiando «y» con «o».
    \end{block}
\end{frame}

\begin{frame}{Ejemplo: De Morgan en un universo finito}
    Sea $X=\Set{1,2,3,4,5,6}$, $A=\Set{1,2,3,4}$ y $B=\Set{3,4,5}$.
    \begin{align*}
        (A\cup B)^c &= \Set{1,2,3,4,5}^c=\Set{6},\\
        A^c\cap B^c &= \Set{5,6}\cap\Set{1,2,6}=\Set{6}.
    \end{align*}
    \begin{exampleblock}{Coinciden}
        $(A\cup B)^c=A^c\cap B^c$.
    \end{exampleblock}
    \centering\small 
\end{frame}

\section{Producto cartesiano}
\begin{frame}{Pares ordenados y producto cartesiano}
    \begin{block}{Par ordenado}
        $(a,b)=(c,d)$ si y solo si $a=c$ y $b=d$. El orden sí importa: en general $(a,b)\ne(b,a)$.
    \end{block}
    \begin{definition}[Producto cartesiano]
        Para conjuntos $A$ y $B$,
        \[
          A\times B=\Set{(a,b):a\in A,\ b\in B}.
        \]
    \end{definition}
    \begin{exampleblock}{Ejemplo}
        Si $A=\Set{1,2}$ y $B=\Set{x,y,z}$, entonces
        \[
          A\times B=\Set{(1,x),(1,y),(1,z),(2,x),(2,y),(2,z)}.
        \]
    \end{exampleblock}
    \begin{exampleblock}{Tarea de programación}
        \href{https://codeforces.com/problemset/problem/268/A}{\textit{Games} (Codeforces)}.
    \end{exampleblock}
\end{frame}

\begin{frame}{Diagrama del producto cartesiano}
    \centering
    \begin{tikzpicture}[x=2.25cm,y=1.05cm]
    \fill[LeafGreen!45] (1,2) rectangle (4,3);
    \fill[LeafGreen!25] (0,0) rectangle (1,2);
    \fill[UdeAGreen] (0,2) rectangle (1,3);
    \fill[CardGreen] (1,0) rectangle (4,2);

    \draw[DeepGreen,thick] (0,0) rectangle (4,3);
    \draw[UdeAGreen] (1,0) -- (1,3);
    \draw[UdeAGreen] (2,0) -- (2,3);
    \draw[UdeAGreen] (3,0) -- (3,3);
    \draw[UdeAGreen] (0,1) -- (4,1);
    \draw[UdeAGreen] (0,2) -- (4,2);

    \node[white,font=\large] at (.5,2.5) {$\times$};
    \node[font=\large\bfseries] at (1.5,2.5) {$x$};
    \node[font=\large\bfseries] at (2.5,2.5) {$y$};
    \node[font=\large\bfseries] at (3.5,2.5) {$z$};
    \node[font=\large\bfseries] at (.5,1.5) {$1$};
    \node[font=\large\bfseries] at (.5,.5) {$2$};
    \node[font=\large] at (1.5,1.5) {$(1,x)$};
    \node[font=\large] at (2.5,1.5) {$(1,y)$};
    \node[font=\large] at (3.5,1.5) {$(1,z)$};
    \node[font=\large] at (1.5,.5) {$(2,x)$};
    \node[font=\large] at (2.5,.5) {$(2,y)$};
    \node[font=\large] at (3.5,.5) {$(2,z)$};

    \node[font=\large\bfseries] at (2.5,3.35) {$B$};
    \node[font=\large\bfseries] at (-.38,1.5) {$A$};
\end{tikzpicture}
    \begin{exampleblock}{Interpretación}
        Las filas corresponden a elementos de $A$, las columnas a elementos de $B$ y cada celda contiene un par de $A\times B$.
    \end{exampleblock}
\end{frame}

\begin{frame}{Cardinalidad del producto cartesiano}
    \begin{theorem}
        Si $A$ y $B$ son finitos, entonces
        \[
          \card{A\times B}=\card{A}\,\card{B}.
        \]
    \end{theorem}
    \begin{block}{Demostración}
        Para cada $a\in A$, los pares que comienzan con $a$ son
        \[
          \Set{(a,b):b\in B},
        \]
        y hay $\card{B}$ de ellos. Los grupos correspondientes a distintos $a$ no se superponen. Como hay $\card{A}$ elecciones de primera coordenada, el total es $\card{A}\card{B}$.
    \end{block}
    \begin{exampleblock}{Tarea de programación}
        \href{https://codeforces.com/problemset/problem/228/A}{\textit{Is your horseshoe on the other hoof?} (Codeforces)}.
    \end{exampleblock}
\end{frame}

\end{document}